\section{Discusión, conclusiones y agenda de investigación}

Las últimas tres diapositivas devuelven los resultados al terreno del método de investigación. El ponente no afirma que el CoT ya esté explicado, ni que la prompt engineering vaya a existir para siempre; lo que enumera son las conclusiones que podían obtenerse de los experimentos en ese momento, las preguntas sobre el mecanismo que aún falta responder y su juicio personal sobre las direcciones en las que pueden participar los estudiantes.

\subsection{¿Desaparecerá la prompt engineering?}

En la sesión de preguntas se distinguen las easy, well-specified tasks de las increasingly hard, underspecified tasks. Los modelos más grandes pueden hacer que un multiple-choice sencillo sea más robust frente a la redacción, pero cuanto más capaz es el modelo, más se le emplea para objetivos complejos; en ese caso, cómo describir las restricciones, aportar evidence y especificar la decomposition y la verification sigue siendo interface design.

\lecturefigure{slide-34.jpg}{Chain-of-thought discussion: la persistencia del prompting y las preguntas abiertas}{Diapositiva 34 del material oficial}

\begin{knowledgebox}{Lectura de la figura: separar las conclusiones conocidas de los mecanismos desconocidos}
La evidencia de la clase respalda que el CoT funciona en varios benchmark, idiomas y modelos, y muestra un umbral de escala; queda sin resolver qué experiencias del pretraining hacen que el modelo aprenda a usar reasoning traces, qué relación guardan esas traces con el cómputo interno y qué tareas se generalizan a partir del few-shot CoT. El valor de la diapositiva de Discussion está precisamente en mantener a la vista los unknowns que hay detrás de los resultados exitosos.
\end{knowledgebox}

\teachervoice{Alguien preguntó si el CoT es solo el compute que aportan los token adicionales. El ponente mencionó que el control con relleno de igual longitud (equal-length filler) no funciona, y sostuvo que el lenguaje es importante para guiar el reasoning; cuando alguien preguntó por dar primero la respuesta y luego la explicación, dijo que por lo general es peor. Estos son negative controls importantes, pero aun así calificó explícitamente ``why CoT works'' como una pregunta abierta.}

\subsection{Conclusiones de toda la clase}

Las diapositivas condensan las dos partes de la clase en tres niveles: las capacidades pueden presentar un umbral en la metric de la tarea conforme crece la escala; el prompting, el fine-tuning y la data quality desplazan la frontera medible; el CoT y la self-consistency muestran que un inference-time method también debe corresponder a la base-model capability.

\lecturefigure{slide-35.jpg}{Conclusiones de la clase: una perspectiva unificada de Scaling, CoT y self-consistency}{Diapositiva 35 del material oficial}

\begin{importantbox}{Los seis puntos más valiosos de toda la clase}
\begin{enumerate}
  \item El smooth loss scaling y el thresholded task performance pueden darse al mismo tiempo.
  \item Juzgar la emergence depende de la model family, el prompt, la metric y el baseline.
  \item Better data, architecture y post-training pueden desplazar el umbral empírico hacia modelos más pequeños.
  \item El CoT aporta decomposition y sequential computation mediante token intermedios, pero no garantiza que la explicación sea fiel.
  \item BBH y el multilingual CoT muestran que la ganancia se extiende entre tareas e idiomas, aunque con una heterogeneidad notable.
  \item La self-consistency convierte las inference samples en accuracy; deben reportarse explícitamente el cost y la error correlation.
\end{enumerate}
\end{importantbox}

\subsection{Direcciones de investigación que propone el ponente}

En enero de 2023, la realidad era que escalar sin más los foundation model correspondía sobre todo a unos cuantos laboratorios industriales. Por eso el ponente sitúa las direcciones accesibles para los estudiantes en better prompting, capability characterization, applications, benchmarks y compute-efficient methods: estos trabajos permiten aprovechar mejor los modelos existentes y, a la vez, poner a prueba los límites de la scaling narrative.

\lecturefigure{slide-36.jpg}{Looking forward: los intereses de investigación personales del ponente}{Diapositiva 36 del material oficial}

\begin{knowledgebox}{Lectura de la figura: qué nivel del problema atiende cada una de las cinco direcciones}
La investigación en Scaling amplía la base capability; el better prompting mejora la elicitation; la characterization construye un mapa de capacidades; el applied work lleva los modelos a workflow reales; los benchmarks ofrecen mediciones que no se saturan rápidamente; y los compute-efficient methods reducen la barrera de entrada del entrenamiento y la inferencia. No se sustituyen entre sí, sino que avanzan en conjunto desde cinco niveles: capacidad, interfaz, medición, aplicación y recursos.
\end{knowledgebox}

\teachervoice{Al cerrar, el ponente dijo que los benchmark existentes se resuelven muy rápido y que se necesitan evaluaciones nuevas y más difíciles; también que conviene investigar métodos más compute-efficient, para que no solo la industria pueda usar modelos más capaces. Este cierre coloca en una misma agenda de investigación los beneficios del scaling y el costo de acceso.}

\subsection{Resumen de la sección}

La forma de la prompt engineering puede cambiar, pero definir objetivos claros, descomponer tareas y verificar resultados no desaparecerá por sí solo. La madurez de la clase está en lo siguiente: los resultados de comportamiento son sólidos y la explicación de los mecanismos sigue siendo cautelosa; escalar modelos merece investigación, pero los datos, la evaluación y la eficiencia también determinan si una capacidad puede descubrirse, reproducirse y desplegarse.
